\scanpage{172}
\mspage{121}
% archive sheet number: 172
3. \uline{Sommes amalgamées} \uline{et intégration}

\diagram{\(\displaystyle\int\!\begin{tikzcd}[ampersand replacement=\&,row sep=0pt,column sep=small]
\& X_1\\
X_0\arrow[ur]\arrow[dr] \&\\
\& X_2
\end{tikzcd}\,.\)}

\uline{3.1.} \uline{Réduction canonique au cas d'un diagramme d'immersions \deleted{\unsure{ouvertes}} fermées}.

Considérons un diagramme

\diagram{\(\text{(\unsure{$*$|$\alpha$})}\qquad\begin{tikzcd}[ampersand replacement=\&,row sep=small,column sep=large]
\& X_1\\
X_0\arrow[ur,"\unsure{\alpha_1|a_1}"]\arrow[dr,"\alpha_2"'] \&\\
\& X_2
\end{tikzcd}\)}
dans $\mathrm{Cat}$. Nous pouvons le regarder comme un contrafoncteur

\[
\unsure{\mathscr{I}|\Phi}^{\circ}\longrightarrow\mathrm{Cat}
\]
où

\diagram{\(\displaystyle\mathscr{I}=\left(\begin{tikzcd}[ampersand replacement=\&,row sep=0pt,column sep=small]
\& 1\arrow[dl]\\
0 \&\\
\& 2\arrow[ul]
\end{tikzcd}\right),\)}
donc comme définissant une catégorie fibrée (scindée) \inserted{$\unsure{\mathscr{X}|\mathcal X}$} sur $\mathscr{I}$. Il est parfois commode de représenter cette catégorie $\mathscr{X}$ par le \lq\lq signe d'intégration\rq\rq

\diagram{\(\displaystyle\int\!\begin{tikzcd}[ampersand replacement=\&,row sep=0pt,column sep=small]
\& X_1\\
X_0\arrow[ur]\arrow[dr] \&\\
\& X_2
\end{tikzcd}\,,\)}
et \unsure{la|le} regarder comme une sorte de \lq\lq somme amalgamée homotopique\rq\rq{} de $X_1$ et $X_2$ sur $X_0$, i.e. comme une \lq\lq limite inductive homotopique\rq\rq{} du diagramme de départ ($*$).
